\section{Design of the experiment}\label{sec:design}

\subsection{How the design evolved}\label{sec:evolution}

The first design (v1, \path{docs/design/parallel-measurement-mode-v1.md}) wanted to measure the path not taken
\emph{inside} a live session: at a decision point, fork the session, run the other path on the same state, record
its cost, and throw the result away. It proposed single-call forks for prepared actions, turn forks for routing,
and a deferred launch with a ``tail correction'' to undo cache effects. A pre-flight review rejected three parts
of it, and the second design (v2, \path{docs/design/parallel-measurement-mode.md}) is what ran:
\begin{itemize}
  \item \textbf{Single-call forks are not a savings unit.} The probe showed that when a call is skipped, the cache
    write it would have made simply moves to the next call (\cref{sec:caching}); a one-call fork only re-derives
    the estimate it was meant to check.
  \item \textbf{Turn forks measure the wrong thing.} A fork that switches model in the middle of someone else's
    history measures a switch from a foreign history, not what the policy costs in its own world.
  \item \textbf{Power is counted in scenarios.} Cost ratios vary far more between scenarios than within one, so
    the unit of replication had to be the scenario (\cref{sec:power}).
\end{itemize}
v2 replaced forks with whole paired sessions, isolated by a per-session nonce rather than by timing, and treats
cache state as an outcome to be measured, not a nuisance to be corrected.

\subsection{Paired sessions from one frozen snapshot}

The unit of comparison is a \term{pair}: one session of a routed arm and one session of the plain host
(the \term{anchor}) on the same scenario, repetition and host. Both start from the \emph{same} frozen copy
of the workspace and receive the \emph{same} scripted user messages, so any difference in their bills
comes from the policy, not from the task.

All arms of one scenario, repetition and host form a \term{wave}. A wave's sessions are launched within
seconds of each other, each in its own terminal, so they share the time of day and the provider's load.
The launch order was fixed, not randomised: measured from the first session of its wave, the anchor started
on average \StaggerAnchor\,s later, the A/A copy \StaggerAA\,s, shipped \StaggerShipped\,s (at most
\StaggerShippedMax\,s), sticky \StaggerSticky\,s and Sonnet \StaggerSonnet\,s. A small systematic effect of
launch order cannot be ruled out (\cref{sec:aa}). Each session carries its own
nonce, so no arm can read another arm's cache (\cref{sec:caching}). A per-request audit checked every
session for cache reads it could not have written itself; \AuditFlagged\ sessions were flagged.

\begin{keyidea}
\textbf{Pairing removes most of the noise.} Two scenarios can differ in cost by a factor of ten; two arms on
the same scenario differ far less. By comparing each routed session only with its own anchor, the
analysis looks at the ratio within a pair, and the big differences between scenarios cancel out.
\end{keyidea}

\subsection{Arms and hosts}

\begin{table}[htbp]
\centering\small
\caption{The five arms. Every arm except the Sonnet control runs on both host models.}
\label{tab:arms}
\begin{tabular*}{\textwidth}{@{\extracolsep{\fill}}l >{\raggedright\arraybackslash}p{0.62\textwidth}@{}}
\toprule
Arm & What it is, and what comparing it with the anchor tells us \\
\midrule
anchor & The plain host model, no routing. The baseline every other arm is compared with. \\
aa & A second, identical anchor on a seeded \AAPct\,\% subsample. Its ratio to the anchor should be 1: its
  spread is the run-to-run noise floor. \\
shipped & The fast-decisions bundle as shipped (orch-default): its judge re-chooses the model (Sonnet or the
  host) at the start of each turn; it never switched in the middle of a turn (\ShippedMidturn\ mid-turn
  switches in \NShippedSessions\ sessions). \\
sticky & Uses the same routing judge, but only once, on the first prompt; the session then stays on the
  chosen model and never switches. Shipped versus sticky isolates the cost of switching. \\
sonnet & Plain Sonnet for the whole session (``just use the cheap model''). It does not depend on the host,
  so one run serves as the control for both hosts. \\
\bottomrule
\end{tabular*}
\end{table}

The two hosts are Fable 5.1, an expensive model (\cref{tab:prices}), and Opus 5.5, a cheaper one with very
cheap cache reads. The sticky judge chose Sonnet for
\NStickySonnetOnly\ of \NStickySessions\ sticky sessions (\StickyCheap\ of \StickyDecisions\ per host,
\StickyCheapPct\,\%) and the host for the rest, so in practice the sticky arm is mostly ``Sonnet for the whole
session''. The judge's own calls are not part of any session's cost: the sticky decision was made once by the
harness, and the shipped arm's per-turn decisions went to a separate judge service (Jev). Both used the bundle's
default judge, Jev; the post-hoc Cloudflare Clef judges of \cref{sec:clef} were never run in the campaign. An earlier benchmark
measured Jev at \$\JevPerMillion\ per million decisions on comparable real states; at one decision per turn
that adds about \JudgeShippedCents\ cents to a shipped session, \JudgeShippedShare\,\% of a plain Fable
session's cost. That estimate comes from another study and is not measured here.

\subsection{Scenarios}

There are \NScenarios\ scenarios in \NFamilies\ families: \FamPolyglot\ coding exercises from a
multi-language benchmark (polyglot), \FamRepos\ bug fixes and features in small public repositories
(repos), \FamMixed\ mixed coding-and-analysis tasks (mixed), and \FamKnowledge\ documentation, explanation
and review tasks (knowledge). Each is a fixed script of \MinTurns\ to \MaxTurns\ user turns (\MeanTurns\ on
average), identical for every arm, with automatic checks after each turn and hidden tests at the end.
Turns are \ShortGapS\,s apart, except that \NLongGapScen\ scenarios contain \LongGapsPerScen\ pauses of
\LongGapMin\ minutes, longer than the cache's \CacheTTLMin-minute lifetime, so expiry is part of the
workload. The full list is in \cref{tab:scenarios}.

\paragraph{The preregistered split.} Before any session ran, the scenarios were split with a fixed random
seed (\SplitSeed) into \NTrain\ \term{training} scenarios and \NTest\ \term{test} scenarios, balanced by
family and by whether a scenario has long pauses. The savings model is fitted on training scenarios only;
every confirmatory claim uses test scenarios only.

\begin{table}[htbp]
\centering\small
\caption{Scenarios by family and preregistered split.}
\label{tab:families}
\begin{tabular*}{0.7\textwidth}{@{\extracolsep{\fill}}l r r r@{}}
\toprule
Family & training & test & total \\
\midrule
polyglot & \FamPolyglotTrain & \FamPolyglotTest & \FamPolyglot \\
repos & \FamReposTrain & \FamReposTest & \FamRepos \\
mixed & \FamMixedTrain & \FamMixedTest & \FamMixed \\
knowledge & \FamKnowledgeTrain & \FamKnowledgeTest & \FamKnowledge \\
\midrule
all & \NTrain & \NTest & \NScenarios \\
\bottomrule
\end{tabular*}
\end{table}

\subsection{Quality}

A routed arm that saves money by doing worse work has not saved anything. Each turn's checks give a
pass or fail; the \term{turn-pass fraction} is the share of a session's turns that passed. The hidden tests
at the end give a final pass or fail. A savings claim is allowed only where quality is
\term{non-inferior}: the 95\,\% lower bound of the mean difference in turn-pass fraction (arm minus anchor)
must be above \NIMargin, that is, at most \NIMarginPts\ percentage points worse.

\subsection{What a session costs}

Every request's tokens were priced with the fixed table in \cref{tab:prices}; the recomputed cost matched
the provider's reported cost on every session (\CostMismatch\ mismatches). One adjustment makes the
comparison fair. Every session's first request writes the shared tool definitions (about \ToolsAllowance\
tokens) to a cold cache, because the nonce stops sessions from sharing. In production those tools would
usually be warm. The primary \term{tools-normalized} cost re-prices that shared prefix as a cache read; it
changes the campaign total by \$\ToolsDeltaTotal, and every result is also reported on the raw basis
(\cref{sec:rawbasis}).

\subsection{Hypotheses and the decision rule}

The hypotheses, their thresholds, the train/test split and the decision rule were fixed in the
preregistration. Each is tested on the test split only:
\begin{itemize}
  \item \textbf{H1 (Fable host):} sticky and shipped cost less than the anchor (upper end of the interval
    below 1).
  \item \textbf{H2 (Opus host):} routing to Sonnet does not save: sticky, shipped and Sonnet all have a
    lower end above \HTwoLower.
  \item \textbf{H3 (both hosts):} sticky costs no more than shipped (upper end of sticky/shipped at most
    \HThreeUpper).
  \item \textbf{Quality:} every savings claim has non-inferior turn-pass fraction.
\end{itemize}
The decision rule: on a host where H1 holds with non-inferior quality, route (preferring sticky if H3
holds); on a host where H2 holds, do not route to Sonnet. A savings model fitted on training data must
also pass a check on the test data: its 90\,\% prediction intervals must cover between \CovLow\ and
\CovHigh\ of test pairs, and the measured test total must fall inside its predicted interval.

\paragraph{What was fixed later.} The \emph{estimator} was not in the preregistration. The confirmatory
script, \path{evals/paired_confirm.py} (commit \code{\ConfirmSha}), was written after the data were
collected, about \ConfirmAfterMin\ minutes after the campaign finished, when exploratory summaries over all
scenarios (the savings-model output in \path{model/}) already existed. It fixed the number of bootstrap
resamples (\Boot) and the seed (\BootSeed), the pair-level quality filter, which cells count as savings
claims, and the verdict rule, including what ``contradicted'' means. These choices are listed with the other
clarifications in \cref{sec:deviations}; the verdicts under the alternative readings are reported alongside.

\paragraph{Ratios and intervals.} The cost of a pair is summarized as a \term{ratio}, arm cost divided by
anchor cost: 0.80 means 20\,\% cheaper. Ratios are averaged as a \term{geometric mean}, which treats
``twice as expensive'' and ``half as expensive'' symmetrically. A \term{95\,\% confidence interval} gives
the range of ratios consistent with the data. It is computed with a \term{scenario-cluster bootstrap}:
resample whole scenarios (then repetitions) with replacement thousands of times, recompute the ratio each
time, and take the middle 95\,\%. Resampling scenarios rather than single sessions keeps the interval
honest about how much the conclusions depend on which tasks were chosen.

\subsection{How big the experiment had to be}\label{sec:power}

Earlier data gave the spread of log cost ratios between scenarios and within a scenario. Because the
between-scenario spread dominates, precision is bought with more \emph{scenarios}, not more repetitions.
\Cref{tab:power} shows the half-widths the design document projected; the campaign used \NScenarios\
scenarios and \NReps\ repetitions.

\begin{table}[htbp]
\centering\small
\caption{Projected precision of the routing ratio, from the design document (before the campaign). The
projections assumed \PowerSA\ and \PowerSB\ scenarios; the confirmatory test split has \NTest, so its
intervals are wider than these.}
\label{tab:power}
\begin{tabular*}{\textwidth}{@{\extracolsep{\fill}}l r r r r@{}}
\toprule
 & \multicolumn{2}{c}{SD of log ratio} & \multicolumn{2}{c}{95\,\% CI half-width} \\
\cmidrule(lr){2-3}\cmidrule(l){4-5}
Contrast (earlier data) & between scenarios & within & \PowerSA\ scenarios & \PowerSB\ scenarios \\
\midrule
Routing, Fable (s1m values) & 0.265 & 0.297 & $\pm$9\,\% & $\pm$11\,\% \\
Routing, Fable (S1 values, worst case) & 0.433 & 0.258 & $\pm$13\,\% & $\pm$16\,\% \\
Routing, Opus (S1 values) & 0.162 & 0.266 & $\pm$7\,\% & $\pm$8\,\% \\
Routing, Opus (s1m values) & 0.027 & 0.211 & $\pm$4\,\% & $\pm$5\,\% \\
\bottomrule
\end{tabular*}

\end{table}

Behind that table is a variance analysis of \PowerPairs\ earlier pairs (\path{step0_variance.json}). For the
shipped router on the Fable host in \CaTurns-turn sessions, the between-scenario standard deviation of the log
ratio was \PowFableFourSB\ against \PowFableFourSW\ within a scenario; on Opus it was \PowOpusFourSB\ against
\PowOpusFourSW. \Cref{fig:power} turns these into the interval one should expect for a given number of scenarios.

\begin{figure}[htbp]
\centering
\pgfplotslegendfromname{powerlegend}\par\smallskip
% Projected 95% CI half-width of the cost ratio vs number of scenarios (2 repetitions), from step0_variance.json.
\begin{tikzpicture}
\begin{axis}[benchmark, scale only axis, width=11cm, height=5.2cm, xmin=10, xmax=100, ymin=0, ymax=40,
  xlabel={Number of scenarios (2 repetitions each)}, ylabel={95\,\% CI half-width (\%)}, xmajorgrids, ymajorgrids,
  title={Precision is bought with scenarios}, title style={font=\small\bfseries},
  cycle list={{okblue, mark=*, solid}, {okblue, mark=o, dashed}, {okorange, mark=square*, solid}, {okorange, mark=square, dashed}},
  every axis plot/.append style={line width=1pt, mark size=1.8pt, mark repeat=2},
  legend to name=powerlegend, legend columns=4, legend cell align=left,
  legend style={/tikz/every even column/.append style={column sep=8pt}}]
\addplot+ table[x=S, y=c0]{generated/data/power-curve.dat};
\addplot+ table[x=S, y=c1]{generated/data/power-curve.dat};
\addplot+ table[x=S, y=c2]{generated/data/power-curve.dat};
\addplot+ table[x=S, y=c3]{generated/data/power-curve.dat};
\addlegendentry{Fable, 1-turn}
\addlegendentry{Fable, 4-turn}
\addlegendentry{Opus, 1-turn}
\addlegendentry{Opus, 4-turn}

\draw[okgray, densely dashed, line width=0.8pt] (axis cs:\NTest,0) -- (axis cs:\NTest,40);
\draw[okgray, densely dashed, line width=0.8pt] (axis cs:\NScenarios,0) -- (axis cs:\NScenarios,40);
\node[font=\scriptsize, anchor=south west, fill=white, inner sep=1pt] at (axis cs:\NTest,34) {test split};
\node[font=\scriptsize, anchor=south west, fill=white, inner sep=1pt] at (axis cs:\NScenarios,34) {all scenarios};
\end{axis}
\end{tikzpicture}

\caption{Projected 95\,\% interval half-width of the cost ratio against the number of scenarios (two repetitions
each), from the earlier data's between- and within-scenario variance.}
\label{fig:power}
\howtoread{Lower is more precise. Each line is one earlier configuration (host, session length). The curves fall
quickly at first and then flatten: going from \NTest\ to \NScenarios\ scenarios roughly halves the uncertainty.
The Fable single-turn line is high because those ratios varied a lot between tasks.}
\end{figure}

\subsection{The pilot}

Two small screening pilots preceded the preregistration; the second ran \PilotSessions\ sessions on
\PilotScen\ scenarios. They confirmed the plumbing (nonce isolation, cache audit, quality checks) and
measured the A/A noise: a standard deviation of the log ratio of \PilotSDFable\ on Fable and \PilotSDOpus\
on Opus. Their cost ratios became the hypotheses above, not findings; they are compared with the final
numbers in \cref{sec:interim}.

\begin{figure}[htbp]
\centering
\pgfplotslegendfromname{pilotlegend}\par\smallskip
% Pilot-2 cost ratios against the final all-split ratios (95% CI) per arm and host.
\begin{tikzpicture}
\begin{axis}[benchmark, scale only axis, width=10.5cm, height=5cm, xmode=log, log ticks with fixed point,
  xtick={0.5,0.6,0.8,1,1.2,1.4,1.6}, xmin=0.45, xmax=1.8, ymin=-0.6, ymax=8.6, ytick={0,1,2,3,5,6,7,8},
  yticklabels from table={generated/data/pilot.dat}{label}, xmajorgrids,
  title={Pilot screen versus the full campaign}, title style={font=\small\bfseries},
  xlabel={Cost ratio against the plain host (log scale)},
  legend to name=pilotlegend, legend columns=2, legend cell align=left,
  legend style={/tikz/every even column/.append style={column sep=10pt}}]
\draw[black!70, line width=0.8pt] (axis cs:1,-0.6) -- (axis cs:1,8.6);
\addplot[only marks, mark=*, mark size=2.6pt, color=okblue, error bars/.cd, x dir=both, x explicit,
  error bar style={line width=1pt}] table[x=final, y=y, x error minus=fm, x error plus=fp]{generated/data/pilot.dat};
\addlegendentry{full campaign, all scenarios (95\,\% CI)}
\addplot[only marks, mark=diamond*, mark size=3.2pt, color=okorange] table[x=pilot, y expr=\thisrow{y}-0.3]{generated/data/pilot.dat};
\addlegendentry{pilot-2 (\PilotScenariosN\ scenarios)}
\end{axis}
\end{tikzpicture}

\caption{The second screening pilot (\PilotPairs\ pairs on \PilotScenariosN\ scenarios) against the full campaign
(all scenarios, 95\,\% intervals), for each arm and host. The first pilot was a smoke test of the harness
(\$\SpendPilotOne) whose rows were not packaged.}
\label{fig:pilot}
\howtoread{Diamonds are the pilot, dots and whiskers the final result. On Fable the diamonds sit inside or near
the whiskers; on Opus the pilot overstated the extra cost. A five-scenario pilot was good for finding the
direction, not the size.}
\end{figure}
